\documentclass[10pt]{article}
\usepackage[margin=0.72in]{geometry}
\usepackage[T1]{fontenc}
\usepackage{lmodern}
\usepackage{amsmath,amssymb,amsthm}
\usepackage{booktabs,longtable,array,tabularx}
\usepackage{microtype}
\usepackage[hidelinks]{hyperref}
\newtheorem{theorem}{Theorem}
\newtheorem{proposition}{Proposition}
\newcommand{\cF}{F}
\newcommand{\cL}{L}
\title{Coalition-Size Influence and Pair Interaction\\
of a Frozen H\'enon Label Predicate}
\author{Anonymous}
\date{}

\begin{document}
\maketitle

\begin{abstract}
We compute the complete first- and second-order discrete derivative atlas of
a monotone Boolean predicate on sixteen named labels.  The predicate is true
for 30400 of 65536 supports.  Every one of the 16 first-order rows is resolved
by coalition size, and every one of the 120 unordered-pair rows retains
positive and negative discrete-Hessian counts, exact uniform Banzhaf
normalization, and exact factorial-weighted Shapley interaction.  The faithful
1920-element label group has seven label orbits and 27 pair orbits; the latter
collapse to ten numerical classes.  An independent block-criterion checker, a
symbolic multilinear kernel, deterministic replay, and 27 hostile mutations
certify the receipt.  The result is finite and combinatorial under the scope
firewall \texttt{NO\_BAD\_EULER\_OR\_ROOT\_NUMBER}.
\end{abstract}

\section{Frozen predicate and contribution}
Let $\cL=\{S_1,\ldots,S_{16}\}$.  Partition nine active direction labels as
\[
 B_1=\{S_1\},\quad B_2=\{S_{16}\},\quad
 B_3=\{S_7,S_{15}\},\quad
 B_4=\{S_3,S_4,S_8,S_{11},S_{12}\}.
\]
For $A\subseteq\cL$, define
\[
 \cF(A)=1
 \quad\Longleftrightarrow\quad
 S_9\in A\ \text{ and }\ \#\{r:A\cap B_r\ne\varnothing\}\geq2.
 \tag{1}
\]
The remaining labels $S_2,S_5,S_6,S_{10},S_{13},S_{14}$ are dummies.
Equivalently, $A$ contains one of 25 minimal generating triples.  Exactly
30400 supports satisfy (1).

Scalar first-order values were available upstream.  The new object here is a
coalition-size-resolved first-order atlas, the complete signed second-order
atlas on all $\binom{16}{2}=120$ pairs, and their faithful label-action
quotient.  Equal numerical values are not used to infer orbit equivalence.

\section{Exact conventions}
For $i\notin A$ and $i,j\notin A$, set
\begin{align*}
 \Delta_i\cF(A)&=\cF(A\cup\{i\})-\cF(A),\\
 \Delta_{ij}\cF(A)&=\cF(A\cup\{i,j\})-\cF(A\cup\{i\})
 -\cF(A\cup\{j\})+\cF(A).
\end{align*}
Define the size-resolved counts
\begin{align*}
 c_i(k)&=\sum_{\substack{A\subseteq\cL\setminus\{i\}\\|A|=k}}
 \Delta_i\cF(A),\\
 p_{ij}(k)&=\#\{A:|A|=k,\ \Delta_{ij}\cF(A)=1\},\\
 n_{ij}(k)&=\#\{A:|A|=k,\ \Delta_{ij}\cF(A)=-1\},\qquad
 d_{ij}(k)=p_{ij}(k)-n_{ij}(k).
\end{align*}
No floating-point convention is involved.  The first-order indices are
\[
 \beta_i=2^{-15}\sum_{k=0}^{15}c_i(k),\qquad
 \phi_i=\sum_{k=0}^{15}c_i(k)\frac{k!(15-k)!}{16!},
 \tag{2}
\]
and the pair indices are
\[
 \beta_{ij}=2^{-14}\sum_{k=0}^{14}d_{ij}(k),\qquad
 I_{ij}=\sum_{k=0}^{14}d_{ij}(k)\frac{k!(14-k)!}{15!}.
 \tag{3}
\]

\section{First-order atlas}
\begin{theorem}[Exact first-order values]
The nonzero rows fall into four numerical classes shown in
Table~\ref{tab:first}.  Six dummy rows vanish.  The raw swing total is 40704,
$\sum_i\beta_i=159/128$, and $\sum_i\phi_i=1$.
\end{theorem}

\begin{table}[h]
\centering
\caption{First-order numerical classes.  Each tuple is (labels per class,
raw swings, uniform Banzhaf, Shapley--Shubik).}
\label{tab:first}
\begin{tabular}{@{}lrrr@{}}
\toprule
Labels & swings & $\beta_i$ & $\phi_i$\\
\midrule
$S_9$ & 30400 & $475/512$ & $271/360$\\
$S_1,S_{16}$ & 2240 & $35/512$ & $61/1260$\\
$S_7,S_{15}$ & 2112 & $33/512$ & $2/45$\\
$S_3,S_4,S_8,S_{11},S_{12}$ & 320 & $5/512$ & $31/2520$\\
$S_2,S_5,S_6,S_{10},S_{13},S_{14}$ & 0 & 0 & 0\\
\bottomrule
\end{tabular}
\end{table}

The complete coalition-size vectors, split after $k=7$, are
\begin{center}
\scriptsize
\begin{tabular}{@{}l p{0.35\textwidth} p{0.35\textwidth}@{}}
\toprule
row & $c_i(0),\ldots,c_i(7)$ & $c_i(8),\ldots,c_i(15)$\\
\midrule
$S_1,S_{16}$ & $0,0,8,59,196,390,521,494$ & $339,166,55,11,1,0,0,0$\\
$S_7,S_{15}$ & $0,0,7,52,175,355,486,473$ & $332,165,55,11,1,0,0,0$\\
$S_3,S_4,S_8,S_{11},S_{12}$ & $0,0,4,25,66,95,80,39$ & $10,1,0,0,0,0,0,0$\\
$S_9$ & $0,0,25,224,940,2461,4504,6095$ & $6269,4950,2992,1364,455,105,15,1$\\
\bottomrule
\end{tabular}
\end{center}

\section{Signed pair interaction atlas}
\begin{theorem}[Ten exact numerical classes]
The 120 unordered pairs occupy the ten classes in
Table~\ref{tab:classes}.  The positive and negative columns give, per pair,
$\sum_kp_{ij}(k)$ and $\sum_kn_{ij}(k)$.  The canonical receipt additionally
retains all 15 cells of $p_{ij}$, $n_{ij}$, and $d_{ij}$ for every pair.
\end{theorem}

\begin{table}[h]
\centering
\caption{Complete numerical classification of the pair atlas.}
\label{tab:classes}
\small
\begin{tabular}{@{}rllrrrr@{}}
\toprule
class & representative & pairs & positive & negative & $\beta_{ij}$ & $I_{ij}$\\
\midrule
K01 & $S_1,S_2$ & 75 & 0 & 0 & 0 & 0\\
K02 & $S_1,S_3$ & 10 & 64 & 256 & $-3/256$ & 0\\
K03 & $S_1,S_7$ & 4 & 64 & 2048 & $-31/256$ & $-5/84$\\
K04 & $S_1,S_9$ & 2 & 2240 & 0 & $35/256$ & $31/168$\\
K05 & $S_1,S_{16}$ & 1 & 64 & 2176 & $-33/256$ & $-11/168$\\
K06 & $S_3,S_4$ & 10 & 0 & 320 & $-5/256$ & $-1/56$\\
K07 & $S_3,S_7$ & 10 & 64 & 128 & $-1/256$ & $1/168$\\
K08 & $S_3,S_9$ & 5 & 320 & 0 & $5/256$ & $5/84$\\
K09 & $S_7,S_9$ & 2 & 2112 & 0 & $33/256$ & $1/6$\\
K10 & $S_7,S_{15}$ & 1 & 0 & 2112 & $-33/256$ & $-13/168$\\
\bottomrule
\end{tabular}
\end{table}

Positive Hessians record complementary label action at a losing coalition;
negative Hessians record redundancy at a winning threshold.  Both signs can
occur for one pair, so retaining only the signed sum would discard structure.
For example, $S_1,S_7$ has 64 positive and 2048 negative coalitions, whereas
$S_7,S_9$ has 2112 positive and no negative coalitions.

\section{Faithful label and pair quotients}
The lifted ambient object has order 11520, but an order-six factor fixes every
label.  The quotient used here is the faithful 1920-element label image.  It
has seven label orbits:
\[
\begin{gathered}
\{S_1,S_{16}\},\ \{S_2\},\ \{S_3,S_4,S_{11},S_{12}\},\
\{S_5,S_6,S_{10},S_{13},S_{14}\},\\
\{S_7,S_{15}\},\ \{S_8\},\ \{S_9\}.
\end{gathered}
\]
Thus $S_8$ and $S_3$ share first-order numbers without sharing a faithful
orbit.  The action has 27 unordered-pair orbits.  Their size spectrum is
\[
 \{1:5,\ 2:7,\ 4:7,\ 5:3,\ 8:1,\ 10:3,\ 20:1\}.
\]
Table~\ref{tab:orbits} gives every orbit representative and its invariant
interaction row.

{\small
\renewcommand{\arraystretch}{0.94}
\begin{longtable}{@{}llrrrrr@{}}
\caption{All 27 faithful unordered-pair orbits.}\label{tab:orbits}\\
\toprule
orbit & representative & size & positive & negative & $\beta_{ij}$ & $I_{ij}$\\
\midrule
\endfirsthead
\toprule
orbit & representative & size & positive & negative & $\beta_{ij}$ & $I_{ij}$\\
\midrule
\endhead
P01 & $S_1,S_2$ & 2 & 0 & 0 & 0 & 0\\
P02 & $S_1,S_3$ & 4 & 64 & 256 & $-3/256$ & 0\\
P03 & $S_1,S_4$ & 4 & 64 & 256 & $-3/256$ & 0\\
P04 & $S_1,S_5$ & 10 & 0 & 0 & 0 & 0\\
P05 & $S_1,S_7$ & 4 & 64 & 2048 & $-31/256$ & $-5/84$\\
P06 & $S_1,S_8$ & 2 & 64 & 256 & $-3/256$ & 0\\
P07 & $S_1,S_9$ & 2 & 2240 & 0 & $35/256$ & $31/168$\\
P08 & $S_1,S_{16}$ & 1 & 64 & 2176 & $-33/256$ & $-11/168$\\
P09 & $S_2,S_3$ & 4 & 0 & 0 & 0 & 0\\
P10 & $S_2,S_5$ & 5 & 0 & 0 & 0 & 0\\
P11 & $S_2,S_7$ & 2 & 0 & 0 & 0 & 0\\
P12 & $S_2,S_8$ & 1 & 0 & 0 & 0 & 0\\
P13 & $S_2,S_9$ & 1 & 0 & 0 & 0 & 0\\
P14 & $S_3,S_4$ & 4 & 0 & 320 & $-5/256$ & $-1/56$\\
P15 & $S_3,S_5$ & 20 & 0 & 0 & 0 & 0\\
P16 & $S_3,S_7$ & 8 & 64 & 128 & $-1/256$ & $1/168$\\
P17 & $S_3,S_8$ & 4 & 0 & 320 & $-5/256$ & $-1/56$\\
P18 & $S_3,S_9$ & 4 & 320 & 0 & $5/256$ & $5/84$\\
P19 & $S_3,S_{11}$ & 2 & 0 & 320 & $-5/256$ & $-1/56$\\
P20 & $S_5,S_6$ & 10 & 0 & 0 & 0 & 0\\
P21 & $S_5,S_7$ & 10 & 0 & 0 & 0 & 0\\
P22 & $S_5,S_8$ & 5 & 0 & 0 & 0 & 0\\
P23 & $S_5,S_9$ & 5 & 0 & 0 & 0 & 0\\
P24 & $S_7,S_8$ & 2 & 64 & 128 & $-1/256$ & $1/168$\\
P25 & $S_7,S_9$ & 2 & 2112 & 0 & $33/256$ & $1/6$\\
P26 & $S_7,S_{15}$ & 1 & 0 & 2112 & $-33/256$ & $-13/168$\\
P27 & $S_8,S_9$ & 1 & 320 & 0 & $5/256$ & $5/84$\\
\bottomrule
\end{longtable}
}

\section{Identities, verification, and scope}
\begin{proposition}[Endpoint identity]
For every $i\in\cL$,
\[
 \sum_{j\ne i}I_{ij}
 =\Delta_i\cF(\cL\setminus\{i\})-\Delta_i\cF(\varnothing).
 \tag{4}
\]
Only $S_9$ has right-hand side one; the other fifteen right-hand sides vanish.
Consequently $\sum_{i<j}I_{ij}=1/2$.  Direct enumeration also gives
$\sum_{i<j}\beta_{ij}=-119/256$; no Banzhaf efficiency claim is made.
\end{proposition}
\begin{proof}
Insert (3), sum over $j\ne i$, and group coalitions by size.  The weighted
sum telescopes from the marginal of $i$ at the empty coalition to its marginal
at the full complementary coalition.  Equation (1) evaluates the endpoints.
\end{proof}

The producer constructs the truth table by testing containment of the 25
minimal triples.  A separate checker reconstructs (1) from the pivot and four
blocks.  A third kernel expands the unique multilinear Boolean polynomial:
it has degree ten and 475 terms, agrees on all 65536 Boolean inputs, and its
symbolic discrete derivatives recover all 136 coalition-size rows.  Clean
replay is byte-identical, and 27 semantic mutations are rejected.

The distance-one autocorrelation from the hash-bound C82 receipt and the
monotone boundary swings partition the sixteen outgoing cube edges from each
true support.  Therefore
\[
 40704+445696=16\cdot30400=486400.
 \tag{5}
\]

The canonical evidence digest is
\[
\begin{gathered}
\texttt{bedeb7a3d912330e5eadc72629ee24d7}\\[-2pt]
\texttt{73648993f73f20f23eaf477028334d6e}.
\end{gathered}
\]
The computation concerns a finite named Boolean predicate and its faithful
permutation action.  It asserts no arithmetic/local data, Euler factor, root
number, automorphy, full Burnside ring/table of marks, or Hilbert--Polya
operator.

\end{document}
